\scanpage{19}
\deleted{$E$}

$A$ courbe elliptique

\[
Y=A\times P^{\deleted{\unsure{k|2|1}}1}\qquad\text{$P$ courbe non. sing. complète conn.}
\]

\[
\begin{aligned}
\Theta_0={}&\Delta-A\times(0)-(0)\times A\quad\text{\deleted{diviseur sur} diviseur can. sur } A\times A\\
&\text{définissant l'autodualité}
\end{aligned}
\]

\[
\Theta=\Theta_0\times P^{\deleted{\unsure{k|1}}1}_{p}\qquad\text{diviseur correspondant sur }A\times A\times P^{\unsure{1|k}}=A\times Y
\]

$c$, classe de diviseur sur $Y$, défini par un \deleted{\unsure{fibré}} faisceau inv. $\underline L$ comme $c=\mathrm{cl}(\underline L)$

\[
X=\overline{V(\underline L)}\quad\text{(clôture projective)}.\quad Y^2\xrightarrow{i}X^3\ \text{section nulle.}
\]

le cycle $\Theta$ sur $A\times Y$ définit

\[
z=(\mathrm{id}_A\times i')_*(\Theta)
\]

cycle de codim $2$ sur $\deleted{\unsure{X}}\,A\times X^3$. On veut calculer, si $s:A\times X^3\unsure{\longrightarrow|-}A$ est la projection, $s_*(z^2)$, et on va donner (pour \deleted{des} choix convenables de $c$) des cas où il est $>0$, d'autres où il est $<0$.

Pour ceci, notons que si $i'=\mathrm{id}_A\times i$, on a

\diagram{\begin{tikzcd}[ampersand replacement=\&,column sep=small]
A\times Y \arrow[dr,"t"] \& \& A\times X \arrow[ll,"i'"'] \arrow[dl,"s"]\\
\& A \&
\end{tikzcd}}

\[
\begin{aligned}
z^2&=i'_*(\Theta)^2=i'_*(\Theta^2c')\\
\text{où}\quad c'&=i'^{*}(i'_*(1))=c^1\ (\text{faisceau normal pour }i')\\
&=1_A\times c,\quad c\ \text{étant le faisceau normal}\\
&\hspace{7em}\text{\unsure{pour|par} }\deleted{\unsure{Y}}\,i:Y\unsure{-|\longrightarrow} X.
\end{aligned}
\]

% archive sheet number: 18
